% =============================================================
% Working with LLMs
% Chapter order is set in ../main.tex; the rendered chapter number
% comes from \chapter{} below.
% =============================================================

\chapter{Working with LLMs}
\label{ch:working}

\todo{Chapter scaffold. The setup-and-wrapper sections are written
in full; the surrounding sections are stubs. Where the material is
operational and dates fast --- model lineups, prices, SDK names ---
the stub also notes how to write it so the chapter ages gracefully.}

Where Chapter~\ref{ch:mental} addressed \emph{what a language model
is}, this chapter addresses \emph{how to call one}. The treatment
is operational. We name the players, give the criteria for choosing
among them, set the reader up to actually run a model in one of
three deployment shapes, define the \code{query\_llm} wrapper that
every later script in the book will import, walk through the
chat-message structure that every modern API exposes, treat
structured outputs as the specialisation of message structure that
powers tool calls, and close with the cost, caching, and latency
considerations that shape every production system. None of this is
theoretically deep; all of it is practically necessary.

\section{Model Providers}
\label{sec:work-providers}

\todo{Short, deliberately. Two parts.}

\paragraph{The landscape.}
\todo{As of 2026, the hosted frontier is dominated by three
providers --- OpenAI (GPT family), Anthropic (Claude family),
Google (Gemini family). The open-weights side has its own gravity:
Meta's Llama, Alibaba's Qwen, DeepSeek, Mistral, Microsoft's Phi,
and Google's Gemma straddling both camps. The point of this
paragraph is to name the players the reader will see referenced
elsewhere in the book; it is not to rank them.}

\paragraph{Axes that age slowly.}
\todo{Rather than a player-by-player capability table, list the
axes on which providers differentiate themselves: hosted vs
open-weights; frontier capability vs commodity-cheap; generalist
vs reasoning-specialised; long context vs strong-but-short;
multimodal vs text-only; permissive vs strict on content. Axes
age much more slowly than rankings, and the reader will use them
to think about whichever providers exist when they are reading.}

\section{Choosing a Model}
\label{sec:work-choose}

\todo{The decision criteria. Capability needed for the task;
context budget; latency tolerance; cost ceiling; whether the data
must stay on premises; whether thinking-mode helps. The
recommended default: use a small, fast, cheap model where you can
and route to a strong one only when you must. Mention orchestration
patterns briefly --- a router model that decides which strong model
sees the task --- and forward-reference the agents chapters where
multi-model orchestration is treated.}

\subsection{Reasoning models in the choice}
\label{sec:work-reasoning}

\todo{Treat ``reasoning'' as one axis of the choice rather than a
separate section. When extended thinking helps (multi-step
mathematics, planning, code that requires careful tracing); when
it does not (lookup, summarisation, simple extraction). Cost
trade-off and latency implications.}

\section{Three Ways to Run a Model}
\label{sec:work-deploy}

The book treats the language model as a service one calls. There
are three deployment shapes for ``service one called'' that any
practitioner will encounter, and the choice between them is mostly
orthogonal to which model is being served. Before continuing, a
reader should have at least one of the three working: any of the
later chapters' scripts will run against any of them, with one
environment variable separating the configurations.

\subsection{The hosted API}
\label{sec:work-deploy-hosted}

The simplest route is to pay one of the major providers per token
and POST to their HTTPS endpoint. OpenAI, Anthropic, and Google
Gemini all expose chat-completion APIs that are similar in shape
and trivially callable from Python with \code{requests}. Setup is
one signup, one API key, one environment variable.

\begin{lstlisting}[language=bash,caption={Hosted-API setup, OpenAI flavour.},label=lst:setup-hosted]
$ pip install requests
$ export OPENAI_API_KEY=sk-...
\end{lstlisting}

A minimal smoke test:

\begin{lstlisting}[language=Python,caption={Smoke test against the OpenAI chat-completion endpoint.},label=lst:smoke-hosted]
import os, requests

rsp = requests.post(
    "https://api.openai.com/v1/chat/completions",
    headers={"Authorization": f"Bearer {os.environ['OPENAI_API_KEY']}"},
    json={
        "model":    "gpt-4o-mini",        # any current chat model
        "messages": [{"role": "user", "content": "say hi in three words"}],
    },
    timeout=60,
)
print(rsp.json()["choices"][0]["message"]["content"])
\end{lstlisting}

\noindent
If this prints something like \emph{``Hi there, friend!''}, the
hosted route is working. The Anthropic and Gemini equivalents
differ only in the URL, header conventions, and one or two field
names; we treat all three uniformly in \S\ref{sec:work-wrapper}.

When this is the right choice: you do not have a GPU, your task
needs a frontier capability, you do not require pinned weights, the
budget is acceptable. The disadvantages were treated in
\S\ref{sec:intro-approach}: hosted models change underneath their
identifiers, and reproducibility against a fixed identifier is not
guaranteed.

\subsection{Local with Ollama}
\label{sec:work-deploy-local}

Run a model on your own machine. The book's default working
configuration is Gemma~4 served via Ollama, for the reproducibility
reasons of \S\ref{sec:intro-approach}. Setup is two steps:

\begin{enumerate}[label=(\arabic*)]
\item Install Ollama. Download from \url{https://ollama.com/download}
  (single binary on macOS, Linux, Windows; or \code{brew install
  ollama} on macOS).
\item Pull a model:

\begin{lstlisting}[language=bash]
$ ollama pull gemma4:e4b      # small, fits on a laptop
$ ollama pull gemma4:26b      # large, needs a real GPU
\end{lstlisting}
\end{enumerate}

\noindent
Ollama listens on \url{http://localhost:11434} by default. A smoke
test:

\begin{lstlisting}[language=Python,caption={Smoke test against a local Ollama server.},label=lst:smoke-local]
import requests

rsp = requests.post(
    "http://localhost:11434/api/chat",
    json={
        "model":    "gemma4:e4b",
        "stream":   False,
        "messages": [{"role": "user", "content": "say hi in three words"}],
    },
    timeout=120,
)
print(rsp.json()["message"]["content"])
\end{lstlisting}

\noindent
When this is the right choice: you have a GPU or Apple Silicon
machine with enough memory, you want pinned-weights
reproducibility, your data is sensitive enough that an external API
is a non-starter, or you prefer a flat hardware cost over metered
per-token billing. The disadvantages: smaller models are noticeably
less capable on the hardest tasks, and the larger Gemma sizes will
not fit on consumer GPUs.

A rough hardware sanity guide for Gemma~4:

\begin{description}[style=nextline,leftmargin=2em,itemsep=0.2em]
\item[\textnormal{Apple Silicon, 16~GB unified memory:}]
  \code{gemma4:e2b} comfortable, \code{gemma4:e4b} usable, anything
  larger swaps.
\item[\textnormal{Apple Silicon, 32~GB unified memory:}]
  \code{gemma4:e4b} comfortable, \code{gemma4:26b} runs but is slow.
\item[\textnormal{NVIDIA GPU, 8~GB VRAM:}] \code{gemma4:e2b} only.
\item[\textnormal{NVIDIA GPU, 24~GB VRAM:}] \code{gemma4:e4b}
  comfortable, \code{gemma4:26b} usable.
\item[\textnormal{NVIDIA GPU, 80~GB VRAM (data-centre class):}]
  everything, including reasoning-mode variants.
\end{description}

\noindent
Numbers are approximate and quantisation-dependent. The reader who
exceeds their hardware by one tier will get correct outputs at
unreasonable speed; the reader who exceeds by two tiers will run
out of memory.

\subsection{Remote model server}
\label{sec:work-deploy-remote}

The third deployment shape is the local model server's pattern,
except the GPU lives in someone else's data centre. You rent a GPU
instance with Ollama (or vLLM, or text-generation-inference)
pre-installed; you get back a URL; the rest of the book's code,
including the \code{query\_llm} wrapper of
\S\ref{sec:work-wrapper}, points at that URL and otherwise behaves
identically to the local case.

This is the right choice when you want pinned open weights but
cannot fit the model on your own hardware. RunPod and Vast charge
by the hour at rates that, for an idle box, are higher than per-token
hosted APIs but for a workload with steady throughput are
competitive. Together and Replicate are slightly different: they
expose hosted open-weight models via APIs (closer in shape to
OpenAI than to Ollama), but the model versions are pinned and the
weights are downloadable elsewhere.

A typical setup with RunPod:

\begin{enumerate}[label=(\arabic*)]
\item Pick a GPU template that has Ollama pre-installed (or use the
  bare Ubuntu template and install it manually).
\item Expose port 11434.
\item Pull the model on first boot
  (\code{ollama pull gemma4:26b}).
\item Point the book's \code{OLLAMA\_BASE\_URL} environment
  variable at the RunPod-issued URL.
\end{enumerate}

\noindent
The smoke test is identical to Listing~\ref{lst:smoke-local} with
\code{localhost:11434} replaced by the public URL. The
\code{query\_llm} wrapper handles the URL substitution
transparently.

\section{Your First Call: a query\_llm Wrapper}
\label{sec:work-wrapper}

The three deployment shapes of \S\ref{sec:work-deploy} expose
subtly different HTTP endpoints. Rather than scatter the
differences across every script in the book, we define a single
\code{query\_llm} function that hides them. Every later chapter
imports this function; the reader can switch between hosted and
local backends with one environment variable.

The wrapper sits at \fname{manuscript/scripts/query\_llm.py}; the
salient pieces are inline below. It is deliberately minimal. There
is no retry logic, no exponential back-off, no token counting, no
streaming, and no tool-call handling. Each of those would be
useful in a production wrapper; none of them belongs in the
function the book uses for thirty pages of teaching. Tool calls are
introduced in Chapter~\ref{ch:toolcalls}, which builds a richer
wrapper on the same backend functions defined here.

\begin{lstlisting}[language=Python,caption={\fname{query\_llm.py} --- a three-backend chat wrapper.},label=lst:query-llm]
import json, os
from typing import Optional

import requests

def query_llm(
    messages: list[dict[str, str]],
    *,
    model:       Optional[str] = None,
    temperature: float = 0.0,
    max_tokens:  int = 2048,
) -> str:
    """Send `messages` to whichever backend LLM_BACKEND selects.

    LLM_BACKEND is one of {"ollama", "openai", "anthropic"};
    defaults to "ollama". Returns the assistant's text response.
    """
    backend = os.environ.get("LLM_BACKEND", "ollama")
    dispatch = {
        "ollama":    _ollama_chat,
        "openai":    _openai_chat,
        "anthropic": _anthropic_chat,
    }
    if backend not in dispatch:
        raise ValueError(f"unknown LLM_BACKEND: {backend!r}")
    return dispatch[backend](messages, model, temperature, max_tokens)


def _ollama_chat(messages, model, temperature, max_tokens) -> str:
    base = os.environ.get("OLLAMA_BASE_URL", "http://localhost:11434")
    rsp = requests.post(
        f"{base}/api/chat",
        json={
            "model":    model or os.environ.get("OLLAMA_MODEL", "gemma4:e4b"),
            "stream":   False,
            "messages": messages,
            "options":  {"temperature": temperature, "num_predict": max_tokens},
        },
        timeout=120,
    )
    rsp.raise_for_status()
    return rsp.json()["message"]["content"]


def _openai_chat(messages, model, temperature, max_tokens) -> str:
    rsp = requests.post(
        "https://api.openai.com/v1/chat/completions",
        headers={"Authorization": f"Bearer {os.environ['OPENAI_API_KEY']}"},
        json={
            "model":       model or "gpt-4o-mini",
            "messages":    messages,
            "temperature": temperature,
            "max_tokens":  max_tokens,
        },
        timeout=120,
    )
    rsp.raise_for_status()
    return rsp.json()["choices"][0]["message"]["content"]


def _anthropic_chat(messages, model, temperature, max_tokens) -> str:
    # Anthropic puts the system message in a separate field rather
    # than in the messages list, so we split it out.
    system = next((m["content"] for m in messages if m["role"] == "system"), "")
    msgs   = [m for m in messages if m["role"] != "system"]
    rsp = requests.post(
        "https://api.anthropic.com/v1/messages",
        headers={
            "x-api-key":         os.environ["ANTHROPIC_API_KEY"],
            "anthropic-version": "2023-06-01",
            "content-type":      "application/json",
        },
        json={
            "model":       model or "claude-haiku-4-5-20251001",
            "system":      system,
            "messages":    msgs,
            "temperature": temperature,
            "max_tokens":  max_tokens,
        },
        timeout=120,
    )
    rsp.raise_for_status()
    return rsp.json()["content"][0]["text"]


if __name__ == "__main__":
    answer = query_llm([
        {"role": "user", "content": "Say hi in three words."},
    ])
    print(answer)
\end{lstlisting}

\noindent
A first call against a local Ollama:

\begin{lstlisting}[language=bash]
$ export LLM_BACKEND=ollama
$ export OLLAMA_MODEL=gemma4:e4b
$ python query_llm.py
Hi there, friend!
\end{lstlisting}

\noindent
Switching to a hosted backend is one variable:

\begin{lstlisting}[language=bash]
$ export LLM_BACKEND=openai
$ export OPENAI_API_KEY=sk-...
$ python query_llm.py
Hello, world there!
\end{lstlisting}

\noindent
Three things about the design.

First, the wrapper accepts and returns plain Python --- a list of
message dicts in, a string out. No SDK objects, no streaming
iterator types, no tool-call envelopes. Per
\S\ref{sec:intro-approach}, the book's chapters compose this
primitive into more interesting behaviour rather than reach for
richer abstractions.

Second, the dispatcher is one dict and one lookup. Switching
backends is an environment-variable change, never a code change.
The reader who has set up only the local backend can ignore the
hosted branches; the reader who has set up only the hosted backend
can ignore the local branch. Neither has to be removed for the
other to work.

Third, the function signature does not include tool calls. Tool
calls add structural complexity (the response can be either a text
reply or a list of tool requests) and naming complexity (which
provider's tool-call schema do we adopt as the unified shape?).
Chapter~\ref{ch:toolcalls} introduces a richer wrapper that handles
tool calls; it builds on the same backend implementations defined
here.

The reader who has run the smoke test of \S\ref{sec:work-deploy}
against at least one backend, and seen \code{query\_llm} return
their first answer, is ready for the rest of the book.

\section{Chats: the Message Structure}
\label{sec:work-chats}

\todo{The four-role abstraction every modern API exposes ---
\code{system}, \code{user}, \code{assistant}, \code{tool} --- and
how those roles map onto the actual token sequence the model
sees. A short worked example using the OpenAI Python SDK and a
matching example with Ollama's chat endpoint to show that the
abstraction is provider-agnostic. Mention where chat templates
diverge under the hood (Gemma~4's special tokens, Claude's XML
conventions) but keep that detail short.}

\subsection{Multi-turn conversations}
\label{sec:work-multiturn}

\todo{The conversation as an append-only list of messages. Why the
host has to re-send the whole list on every turn (the model is
stateless). The implication for cost: long conversations grow
input cost linearly per turn.}

\subsection{The system prompt in practice}
\label{sec:work-system}

\todo{What goes in a system prompt vs what goes in the user
message. Stable instructions go in system; per-task content goes
in user. Forward-reference Chapter~\ref{ch:tools} for
prompting craft proper.}

\section{Structured Outputs}
\label{sec:structured}

By default a language model emits a stream of tokens that decode
to free-form natural language. Many downstream uses need more: a
JSON object with a fixed shape, an XML document, a SQL query, a
function signature. Three families of techniques are in common
use, in increasing order of guarantees:

\begin{enumerate}[label=(\arabic*)]
\item \textbf{Prompting.} The system prompt instructs the model
  to emit, say, ``a single JSON object on one line, no prose, no
  code fences''. This works most of the time on a sufficiently
  large model and almost always with one or two few-shot
  examples. It is the cheapest option and the one most readers
  will start with.
\item \textbf{Constrained decoding.} The runtime intercepts the
  token stream and masks out any token that would violate a
  target grammar (\code{llama.cpp} grammars, \code{Outlines},
  \code{lm-format-enforcer}, the structured-output modes shipped
  by hosted providers). The output is \emph{guaranteed} to parse
  against the schema.
\item \textbf{Fine-tuning.} The model is trained on examples of
  the desired structured format. Used together with constrained
  decoding it produces the most reliable output, at the cost of
  a training run.
\end{enumerate}

\noindent
A tool call is the canonical application of structured outputs: a
JSON object with a known schema, intercepted by the runtime,
dispatched, and answered. Chapter~\ref{ch:toolcalls} treats the
tool call as the unit of work and assumes structured-output
reliability is a solved problem at the layer below.

\todo{Expand: a worked example of constrained decoding against a
local model; a comparison of the schema languages used by major
providers (OpenAI's \code{response\_format}, Anthropic's tool-use
schema, Gemma~4's auto-derived schema from Python type hints);
the trade-off between prompt-engineering reliability and
constrained-decoding rigidity.}

\section{Cost, Caching, and Latency}
\label{sec:work-cost}

\todo{The operational triple. Three subsections.}

\subsection{Cost}
\label{sec:work-cost-cost}

\todo{Per-token pricing. The asymmetry between input and output
tokens (input is much cheaper, often by a factor of three to
five). How to estimate a session cost from a prompt template and
expected output length. The trap of agentic loops: cost can scale
with the number of iterations, not with the user-perceived task
size. Forward-reference the context discussion in
Chapter~\ref{ch:toolcalls}.}

\subsection{Prompt and KV caching}
\label{sec:work-cost-cache}

\todo{Why caching matters disproportionately for long-running
agents (the same system prompt and tool catalogue is sent every
turn). The provider-side surface: explicit cache APIs (Anthropic),
automatic prefix caching (OpenAI, several open-source servers),
nothing-at-all. How to design a system prompt to be
cache-friendly: stable prefix, dynamic suffix.}

\subsection{Latency}
\label{sec:work-cost-latency}

\todo{Time-to-first-token vs total-tokens-per-second; the two
matter for different use cases (interactive UX cares about TTFT,
batch agents care about total throughput). Streaming as a UX
pattern. The latency tax of reasoning models (the thinking
tokens are still tokens). Practical numbers will date; cite the
ratios rather than absolute figures.}
